%%%%%%%%%%%%%%%%%%%%%%%%%%%%%%%%%%%%%%%%%%%%%%%%%%%%%%%%%%%%%%%%%%%%%%
% Lớp 11 | Giới hạn. Hàm số liên tục | Hàm số liên tục
% Dạng 1D3?3-5: Bài toán phương trình có nghiệm
% Số câu: 51  (sinh tự động bởi tach_cau.py; từ nay thầy sửa trực tiếp file này)
%%%%%%%%%%%%%%%%%%%%%%%%%%%%%%%%%%%%%%%%%%%%%%%%%%%%%%%%%%%%%%%%%%%%%%

%%%%%%%%%%%%%%%% TRẮC NGHIỆM 4 PHƯƠNG ÁN %%%%%%%%%%%%%%%%

% nguồn: Ok_Toan11_TracNghiem_Lan2
\begin{ex}%[1D3H3-5]%[ma:0dcf0d0c]
	 Phương trình nào dưới đây có nghiệm trong khoảng $\left(0; 1\right)$
	 \choice
		 {$2x^2-3x+4=0$}
		{$\left(x-1\right)^5-x^7-2=0$} 
		 {$3x^4-4x^2+5=0$}
		 {\True $3x^{2017}-8x+4=0$}
	\loigiai{ 
		Xét hàm  số $f(x)=3x^{2017}-8x+4$.
		Hàm số liên tục  trên đoạn $\left[0; 1\right]$ và $f(0)\cdot f(1)=4\cdot (-1)=-4<0$.
		Vậy phương trình $3x^{2017}-8x+4=0$  có nghiệm trong khoảng $\left(0; 1\right)$.
	} 
\end{ex}

% nguồn: Ok_Toan11_TracNghiem_Lan2
\begin{ex}%[1D3H3-5]%[ma:a147c1fb]
	Cho phương trình $2x^4 - 5x^2 + x + 1 = 0 \quad (1)$. Chọn khẳng định đúng trong các khẳng  định sau:
	 \choice
		{Phương trình $(1)$ có đúng một nghiệm trên khoảng  $(0; 2)$}
		{Phương trình $(1)$  vô nghiệm trên khoảng $(1; 2)$}
		{\True Phương trình $(1)$  có ít nhất hai nghiệm trên khoảng $(0; 2)$}
		{Phương trình $(1)$ vô nghiệm  trên khoảng $(0; 1)$}
	\loigiai{ 
		Xét hàm số $f(x)  = 2x^4 - 5x^2 + x + 1$ liên tục trên $\mathbb{R}$.\\
		Suy ra hàm số liên tục  trên $[-2; 3]$.\\
		Ta  có
		 \[
		\heva{ 
		 &f(0) = 1 \\
		&f(1)  = -1
		} \Rightarrow  f(0) \cdot f(1) < 0.
		 \]
		 Suy ra $f(x) = 0$ có ít nhất một nghiệm thuộc $(0; 1)$.
		\[ 
		\heva{ 
		&f(1) = -1  \\
		 &f(2) = 15
		} \Rightarrow f(1) \cdot f(2) < 0. 
		\] 
		Suy ra $f(x) = 0$ có ít nhất một  nghiệm thuộc $(1; 2)$.\\
		Do  đó, phương trình có ít nhất hai nghiệm thuộc khoảng $(0; 2)$.
	} 
\end{ex}

% nguồn: Ok_Toan11_TracNghiem_Lan2
\begin{ex}%[1D3H3-5]%[ma:25e539df]
	Cho phương trình $4x^4+2x^2-x-3=0$ $(1)$.  Mệnh đề nào dưới đây đúng?
	\choice 
		{Phương trình $(1)$ vô nghiệm  trên khoảng $\left(-1; 1\right)$}
		{Phương trình $(1)$ có đúng một nghiệm trên  khoảng $\left(-1; 1\right)$}
		{Phương trình $(1)$ có  đúng hai nghiệm trên khoảng $\left(-1; 1\right)$}
		{\True Phương trình $(1)$ có ít  nhất hai nghiệm trên khoảng $\left(-1; 1\right)$}
	\loigiai{ 
		Xét $f(x)=4x^4+2x^2-x-3=0$ trên khoảng $\left[-1; 1\right]$. 
		Ta có $f(x)$ liên tục trên đoạn $\left[-1; 1\right]$, và $\heva{f\left(-1\right)&=4\\ f(0)&=-3\\ f(1)&=2}\Rightarrow \heva{f\left(-1\right)\cdot f(0) &< 0\\ f(1)\cdot f(0)  &< 0.}$
		Do đó $f(x)=0$  có ít nhất hai nghiệm trên khoảng $\left(-1; 1\right)$.
	 }
\end{ex}

% nguồn: Ok_Toan11_TracNghiem_Lan2
\begin{ex}%[1D3H3-5]%[ma:7b98e1c5]
	Cho phương trình $2x^3-8x-1=0(1)$.  Khẳng định nào sai?
	\choice 
		{Phương  trình không có nghiệm lớn hơn $3$}
		{Phương trình có đúng $3$ nghiệm phân biệt} 
		{\True Phương trình có $2$ nghiệm lớn  hơn $2$}
		{Phương  trình có nghiệm trong khoảng $\left(-5;-1\right)$}
	 \loigiai{
		Hàm số $f(x)=2x^3-8x-1$ liên tục  trên $\mathbb{R}$.
		Do $f\left(-5\right)=-211$, $f\left(-1\right)=5> 0$, $f(2)=-1< 0$, $f(3)=29> 0$ nên phương trình có ít nhất $3$ nghiệm trên $\left(-5;-1\right)$,  $\left(-1;2\right)$, $\left(2;3\right)$.
		Mà phương trình bậc ba có tối đa $3$ nghiệm nên phương trình  có đúng $3$ nghiệm trên $\mathbb{R}$. Do đó C sai.
	} 
\end{ex}

% nguồn: Ok_Toan11_TracNghiem_Lan2
\begin{ex}%[1D3H3-5]%[ma:22f0b64c]
	Cho hàm số $f(x)=x^5-4x^4-4x^3+14x^2+4x-10$. Số  nghiệm của phương trình $f(x)=0$ trên $\mathbb{R}$ là:
	\choice 
		 {$1$}
		 {$3$}
		 {$4$}
		 {\True $5$}
	\loigiai{ 
		Hàm số $f(x)=x^5-4x^4-4x^3+14x^2+4x-10$ là hàm đa thức có tập xác định  là $\mathbb{R}$ nên liên tục trên $\mathbb{R}$.
		Do đó  hàm số liên tục trên mỗi khoảng $\left(-2;-\dfrac{3}{2} \right)$, $\left(-\dfrac{3}{2};0\right)$, $\left(0;1\right)$, $\left(1;2\right)$, $\left(2;5\right)$.
		Ta có $\heva{& f\left(-2\right)=-26  \\ & f\left(-\dfrac{3}{2} \right)=\dfrac{37}{32}} \Rightarrow f\left(-2\right)\cdot f\left(-\dfrac{3}{2} \right) < 0$.
		Suy ra phương trình có  nghiệm thuộc $\left(-2;-\dfrac{3}{2} \right)$.
		Ta  có $\heva{&{f\left(-\dfrac{3}{2} \right)=\dfrac{37}{32}} \\
		 &{f(0)=-10}
		} \Rightarrow f\left(-\dfrac{3}{2}  \right)f(0) < 0$.
		Suy ra phương trình có  nghiệm thuộc $\left(-\dfrac{3}{2};0\right)$.
		Ta có $\heva{&{f(0)=-10} \\ 
		&{f(1)=1} 
		}  \Rightarrow f(0)f(1) < 0$.
		Suy ra phương trình  có nghiệm thuộc $\left(0;1\right)$.
		Ta có $\heva{&{f(1)=1} \\ 
		 &{f(2)=-10}
		} \Rightarrow f(1)f(2) < 0$. 
		 Suy ra phương trình có nghiệm thuộc $\left(1;2\right)$.
		Ta  có $\heva{&{f(2)=-10} \\
		&{f(5)=485} 
		} \Rightarrow f(2)f(5) < 0$. Suy ra phương trình có nghiệm  thuộc $\left(2;5\right)$.
		Như vậy phương trình có ít nhất $5$ nghiệm thuộc khoảng $\left(-2;5\right)$. 
		Tuy nhiên $f(x)=0$ là phương trình bậc  $5$ có nhiều nhất $5$ nghiệm. Vậy phương trình $f(x)=0$ có đúng $5$ nghiệm trên $\mathbb{R}$.
	} 
\end{ex}

% nguồn: Ok_Toan11_TracNghiem_Lan2
\begin{ex}%[1D3H3-5]%[ma:5785ab2b]
	Phương trình $3x^5 + 5x^3 + 10 = 0$ có nghiệm  thuộc khoảng nào sau đây?
	\choice 
		{\True $(-2; -1)$} 
		{$(-10; -2)$} 
		 {$(0; 1)$}
		{$(-1; 0)$} 
	 \loigiai{
		Xét hàm số $f(x) = 3x^5  + 5x^3 + 10$ liên tục trên $\mathbb{R}$.\\
		Suy ra $f(x)$  liên tục trên $[-2; -1]$.\\
		Ta  có
		\[ 
		\heva{ 
		 &f(-2) = -126 \\
		&f(-1) = 2. 
		} 
		 \]
		Suy ra $f(-2) \cdot f(-1) <  0$.\\
		 Do đó $f(x) = 0$ có nghiệm thuộc khoảng $(-2; -1)$.
	 }
\end{ex}

% nguồn: Ok_Toan11_TracNghiem_Lan2
\begin{ex}%[1D3H3-5]%[ma:4cd776e4]
	Phương trình $3x^5+5x^3+10=0$ có nghiệm thuộc  khoảng nào sau đây?
	 \choice
		{\True $\left(-2;-1\right)$} 
		{$\left(-10;-2\right)$} 
		{$\left(0;1\right)$} 
		{$\left(-1;0\right)$} 
	\loigiai{ 
		Đặt $f(x)=3x^5+5x^3+10$. 
		Hàm số $f(x)$ liên  tục trên $\mathbb{R}$ nên $f(x)$ liên tục trên $\left[-2;-1\right]$. (1)
		Ta  có: $\heva{&f\left(-2\right)=-126\\ &f\left(-1\right)=2}$.
		 Suy ra $f\left(-2\right)\cdot f\left(-1\right)=-126\cdot2=-252< 0$, $(2)$
		Từ $(1)$  và $(2)$ suy ra $f(x)=0$ có nghiệm thuộc khoảng $\left(-2;-1\right)$.
	} 
\end{ex}

% nguồn: Ok_Toan11_TracNghiem_Lan2
\begin{ex}%[1D3H3-5]%[ma:bf93cf0a]
	Cho phương trình $2x^4-5x^2+x+1=0(1)$. Chọn khẳng định  đúng trong các khẳng định sau
	 \choice
		{Phương  trình $(1)$ có đúng một nghiệm trên khoảng $\left(-2; 1\right)$}
		{Phương trình $(1)$ vô nghiệm} 
		{\True Phương trình $(1)$ có ít  nhất hai nghiệm trên khoảng $\left(0; 2\right)$}
		{Phương  trình $(1)$ vô nghiệm trên khoảng $\left(-1; 1\right)$}
	 \loigiai{
		Ta có: $\heva{&  f(-1)=-3 \\& f(0)=1 \\ &{f(1)=-1} \\& f(2)=15.}$
	 }
\end{ex}

% nguồn: Ok_Toan11_TracNghiem_Lan2
\begin{ex}%[1D3H3-5]%[ma:792fe683]
	Cho phương trình $2x^3-8x-1=0 \quad (1)$. Khẳng định nào  sai?
	\choice 
		 {Phương trình $(1)$ không có nghiệm lớn hơn $3$}
		{Phương trình $(1)$ có đúng $3$ nghiệm phân biệt} 
		{\True Phương trình $(1)$ có  $2$ nghiệm lớn hơn $2$}
		{Phương trình $(1)$ có nghiệm trong khoảng $(-5; -1)$} 
	 \loigiai{
		Xét hàm số $f(x) =  2x^3 - 8x - 1$ liên tục trên $\mathbb{R}$.\\
		Do đó, hàm số liên tục trên các khoảng $(-5; -1)$, $(-1; 2)$,  và $(2; 3)$.\\
		 Ta có
		 \[
		\heva{ 
		 &f(-5) = -211 \\
		 &f(-1) = 5
		} \Rightarrow  f(-5) \cdot f(-1) < 0.
		\] 
		Suy ra phương trình $f(x) = 0$ có  ít nhất một nghiệm thuộc $(-5; -1)$.
		 \[
		 \heva{
		&f(-1) = 5 \\ 
		&f(2) = -1 
		} \Rightarrow  f(-1) \cdot f(2) < 0.
		 \]
		Suy ra phương trình $f(x)  = 0$ có ít nhất một nghiệm thuộc $(-1; 2)$.
		 \[
		 \heva{
		 &f(2) = -1 \\
		 &f(3) = 29
		} \Rightarrow f(2) \cdot f(3) <  0.
		\] 
		Suy  ra phương trình $f(x) = 0$ có ít nhất một nghiệm thuộc $(2; 3)$.\\
		Như vậy, phương trình $(1)$ có ít ba nghiệm phân biệt trên các khoảng $(-5; -1)$, $(-1;  2)$, và $(2; 3)$.\\
		Mà phương trình $(1)$ là phương trình bậc ba nên phương trình $(1)$ có đúng  ba nghiệm phân biệt trên các khoảng $(-5; -1)$, $(-1; 2)$, và $(2; 3)$.\\
		Vì vậy,  phương trình $(1)$ có đúng $1$ nghiệm lớn hơn $2$.
	 }
\end{ex}

% nguồn: Ok_Toan11_TracNghiem_Lan2
\begin{ex}%[1D3H3-5]%[ma:89849a4f]
	Số nghiệm của phương  trình $x^3-3x-1=0$ là
	 \choice
		{$0$} 
		{\True $1$} 
		{$2$} 
		 {$3$}
	\loigiai{ 
		 Đặt $f(x)=x^3-3x-1$.\\
		Do $f(x)=x^3-3x-1$ là hàm đa thức có tập xác định là $\mathbb{R}$  nên liên tục trên $\mathbb{R}$.\\
		Do  đó hàm số liên tục trên mỗi khoảng $(-2;-1)$, $(-1;0)$, $(0;2)$.\\
		Ta có 
		\begin{itemize}
			\item $\heva{& f(-2)=-3 \\ & f(-1)=1} \Rightarrow f(-2)f(-1)<0 \Rightarrow$ phương trình đã cho có ít  nhất một nghiệm thuộc $(-2;-1).$
			\item $\heva{& f(-1)=1 \\ & f(0)=-1} \Rightarrow f(-1)f(0)<0 \Rightarrow$ phương trình đã cho có  ít nhất một nghiệm thuộc $(-1;0).$
			\item $\heva{& f(2)=1 \\ & f(0)=-1} \Rightarrow f(2)f(0)<0 \Rightarrow$ phương trình đã cho có ít nhất  một nghiệm thuộc $(0;2).$
		\end{itemize}
		Như vậy phương trình đã cho  có ít nhất ba thuộc khoảng $(-2;2)$.\\
		Tuy nhiên phương trình $f(x)=0$ là phương trình bậc ba có nhiều nhất  ba nghiệm.\\
		Vậy phương trình $f(x)=0$ có đúng $3$  nghiệm trên $\mathbb{R}$.
	} 
\end{ex}

% nguồn: Ok_Toan11_TracNghiem_Lan2
\begin{ex}%[1D3H3-5]%[ma:44a673d8]
	Cho phương trình $x^3 - 3x^2 + 2 = 0  \quad (1)$. Chọn khẳng định đúng trong các khẳng định sau:
	 \choice
		{\True Phương trình  $(1)$ có ít nhất hai nghiệm trên khoảng $(-2; 3)$}
		{Phương trình $(1)$  có đúng một nghiệm trên khoảng $(-2; 3)$}
		{Phương trình $(1)$  vô nghiệm trên khoảng $(0; 3)$}
		{Phương trình  $(1)$ vô nghiệm trên khoảng $(-2; 0)$}
	 \loigiai{
		Xét hàm số $f(x) = x^3  - 3x^2 + 2$ liên tục trên $\mathbb{R}$.\\
		Suy ra hàm số liên  tục trên $[-2; 3]$.\\
		Ta  có
		 \[
		\heva{ 
		&f(-2) = -18 \\ 
		&f(0)  = 2
		}  \Rightarrow f(-2) \cdot f(0) < 0.
		\] 
		Suy ra $f(x) = 0$ có ít  nhất một nghiệm thuộc $(-2; 0)$.
		\[ 
		\heva{ 
		&f(0) =  2 \\
		&f(2)  = -2
		} \Rightarrow f(0)  \cdot f(2) < 0.
		 \]
		Suy  ra $f(x) = 0$ có ít nhất một nghiệm thuộc $(0; 2)$.\\
		Do đó, phương trình có ít nhất hai nghiệm thuộc khoảng  $(-2; 3)$.
	 }
\end{ex}

% nguồn: Ok_Toan11_TracNghiem_Lan2
\begin{ex}%[1D3H3-5]%[ma:8b393309]
	Cho phương trình $x^3-3x^2+2=0(1)$. Chọn khẳng định đúng trong các  khẳng định sau?
	\choice 
		{\True Phương trình $(1)$ có ít nhất hai nghiệm  trên khoảng $\left(-2; 3\right)$}
		{Phương trình $(1)$ có đúng một nghiệm  trên khoảng $\left(-2; 3\right)$}
		{Phương trình $(1)$ vô nghiệm} 
		{Phương trình $(1)$ có hai nghiệm trên  khoảng $\left(-2; 0\right)$}
	\loigiai{ 
		Ta có  hàm số $f(x)=x^3-3x^2+2=0$ liên tục trên $\mathbb{R}$.
		Suy ra hàm số liên  trên $\left[-2;3\right]$.
		Ta có: $\heva{& f\left(-2\right)=-18 \\ &f(0)=2} \Rightarrow f(x)=0$ có ít nhất một  nghiệm thuộc $\left(-2;0\right)$.
		 Ta có: $\heva{& f(0)=2 \\ & f(2)=-2} \Rightarrow f(x)=0$ có ít nhất một nghiệm thuộc $\left(0;2\right)$.
		Do đó phương trình có ít nhất  hai nghiệm thuộc khoảng $\left(-2; 3\right)$.
	} 
\end{ex}

% nguồn: Ok_Toan11_TracNghiem_Lan2
\begin{ex}%[1D3H3-5]%[ma:c9231d21]
	Cho hàm số $y = f(x)$ liên tục trên đoạn $\left[a; b\right]$ và thỏa mãn $f(a) = b$, $f(b) = a$ với $a, b >  0$, $a \neq b$. Khi đó phương trình nào sau đây có nghiệm trên khoảng $\left(a; b\right)$?
	 \choice
		{Phương trình $f(x) = 0$} 
		 {\True Phương trình $f(x) = x$}
		{Phương trình $f(x)  = -x$}
		{Phương trình $f(x)  = a$}
	\loigiai{ 
		Xét hàm số  $y = f(x) - x$, hàm số này liên tục trên đoạn $[a; b]$.\\
		 Ta có
		 \[
		 \heva{
		&f(a) -  a = b - a \\
		&f(b) - b = a -  b
		} \Rightarrow \left[f(a) - a\right] \cdot \left[f(b) - b\right] = (b - a)(a  - b) = -(a - b)^2 < 0.
		\] 
		Suy  ra phương trình $f(x) = x$ có nghiệm trên khoảng $(a; b)$.
	} 
\end{ex}

% nguồn: Ok_Toan11_TracNghiem_Lan2
\begin{ex}%[1D3H3-5]%[ma:79292092]
	Cho phương trình $4x^4 + 2x^2  - x - 3 = 0 \quad (1)$. Mệnh đề nào dưới đây đúng?
	 \choice
		 {Phương trình $(1)$ vô nghiệm trên khoảng $(-1; 0)$}
		{Phương trình $(1)$ có đúng một nghiệm trên khoảng  $(-1; 1)$}
		{Phương trình $(1)$ vô nghiệm trên khoảng $(0;  1)$}
		{\True Phương trình $(1)$ có ít nhất hai nghiệm  trên khoảng $(-1; 1)$}
	\loigiai{ 
		Xét hàm số $f(x) = 4x^4 + 2x^2 - x - 3$ trên  khoảng $[-1; 1]$.\\
		Ta có $f(x)$ liên tục  trên đoạn $[-1; 1]$.\\
		Ta  có
		\[ 
		\heva{ 
		&f(-1) =  4 \\
		 &f(0) = -3
		} \Rightarrow f(-1)  \cdot f(0) < 0.
		 \]
		Suy ra $f(x) = 0$ có ít nhất một nghiệm  thuộc $(-1; 0)$.
		\[ 
		 \heva{
		&f(0) = -3  \\
		&f(1) = 2 
		}  \Rightarrow f(0) \cdot f(1) < 0.
		\] 
		Suy  ra $f(x) = 0$ có ít nhất một nghiệm thuộc $(0; 1)$.\\
		Do đó, phương trình có ít nhất hai nghiệm thuộc khoảng $(-1;  1)$.
	 }
\end{ex}

% nguồn: Ok_Toan11_TracNghiem_Lan2
\begin{ex}%[1D3H3-5]%[ma:23392f34]
	Cho hàm số $f(x)$ liên tục trên đoạn $[-1;4]$ sao cho $f(-1)=2$, $f(4)=7$. Có  thể nói gì về số nghiệm của phương trình $f(x)=5$ trên đoạn $[-1;4]$?
	 \choice
		{Vô  nghiệm}
		{\True Có ít nhất một  nghiệm}
		{Có đúng một nghiệm} 
		{Có đúng hai  nghiệm}
	\loigiai{ 
		Ta có $f(x)=5  \Leftrightarrow f(x)-5=0$. Đặt $g(x)=f(x)-5.$\\
		Khi đó $\heva{& g(-1)=f(-1)-5=2-5=-3 \\  & g(4)=f(4)-5=7-5=2} \Rightarrow g(-1)g(4)<0.$\\
		Vậy phương trình $g(x)=0$ có ít nhất một nghiệm thuộc khoảng $(1;4)$ hay phương trình $f(x)=5$ có ít nhất một nghiệm  thuộc khoảng $(1;4)$.
	} 
\end{ex}

% nguồn: Ok_Toan11_TracNghiem_Lan2
\begin{ex}%[1D3V3-5]%[ma:db82882b]
	Cho hàm số $y=f(x)$ liên tục trên đoạn $\left[a;b\right]$ và thỏa mãn $f(a)=b$, $f(b)=a$ với $a,b > 0$, $a\ne b$. Khi đó phương trình nào sau đây có nghiệm  trên khoảng $\left(a;b\right)$.
	 \choice
		{$f(x)=0$} 
		{\True  $f(x)=x$}
		{$f(x)=-x$} 
		 {$f(x)=a$}
	\loigiai{ 
		Hàm số $y=f(x)-x$ liên tục  trên đoạn $\left[a;b\right]$.
		\[\left[f(a)-a\right]\left[f(b)-b\right]=\left(b-a\right)\left(a-b\right)=-\left(a-b\right)^2<  0.\]
		Suy ra  phương trình $f(x)=x$ có nghiệm trên khoảng $\left(a;b\right)$.
	} 
\end{ex}

% nguồn: Ok_Toan11_TracNghiem_Lan2
\begin{ex}%[1D3V3-5]%[ma:79d3ab8a]
	Phương trình nào dưới đây có nghiệm trong  khoảng $(0; 1)$
	 \choice
		{$2x^2  - 3x + 4 = 0$}
		{$3x^4 - 4x^2  + 5 = 0$}
		{$(x-1)^5 -  x^7 - 2 = 0$}
		{\True $3x^{2017} - 8x +  4 = 0$}
	 \loigiai{
		\begin{itemize}
			\item $2x^2 - 3x + 4 = 0$: Phương  trình vô nghiệm.
			\item $3x^4 - 4x^2 + 5 = 0$: Phương trình vô  nghiệm.
			\item Xét phương trình $(x-1)^5 - x^7 -  2 = 0$:\\
			Với $x\in (0;1)$ thì $\heva{&(x-1)^5 \le 0\\&-x^7 \le 0}$ nên $(x-1)^5 - x^7 - 2  < 0$.\\
			Hay phương  trình $(x-1)^5 - x^7 - 2 = 0$ vô nghiệm trong khoảng $(0;1)$.
			\item Xét hàm số  $f(x) = 3x^{2017} - 8x + 4$.\\
			Hàm số liên tục  trên đoạn $[0; 1]$ và ta có
			\[ 
			 \heva{
			&f(0) = 4 \\ 
			 &f(1) = -4
			} \Rightarrow f(0) \cdot f(1) < 0. 
			\] 
			Vậy phương trình $3x^{2017} - 8x + 4 = 0$ có nghiệm trong khoảng $(0; 1)$. 
		\end{itemize}
	 }
\end{ex}

% nguồn: Ok_TracNghiemLop11_Lan1
\begin{ex}%[1D3C3-5]%[ma:9ca7cc41]
	Có tất cả bao nhiêu giá trị nguyên của tham số  $m$ thuộc khoảng $(-10;10)$ để phương trình $x^3-3x^2+(2m+2)x+m-3=0$ có ba nghiệm phân biệt $x_1,x_2,x_3$ thỏa mãn $x_1<-1<x_2<x_3$?
	\choice 
		 {$19$}
		{$18$} 
		{\True  $4$}
		{$3$} 
	 \loigiai{
		Xét hàm số  $f(x)=x^3-3x^2+(2m+2)x+m-3$ liên tục trên $\mathbb{R}$.\\
		Giả sử phương trình có ba nghiệm phân biệt $x_1,x_2,x_3$ sao cho  $x_1<-1<x_2<x_3$. Khi đó $f(x)=(x-x_1)(x-x_2)(x-x_3)$.\\
		 Ta có $f(-1)=(-1-x_1)(-1-x_2)(-1-x_3)>0$ (do $x_1<-1<x_2<x_3$).\\
		Mà $f(-1)=-m-5$ nên suy ra $-m-5>0\Leftrightarrow m<-5$.\\ 
		 Thử lại với $m<-5$, ta có \\
		$\lim\limits_{x\to  -\infty}f(x)=-\infty$ nên tồn tại $a<-1$ sao cho $f(a)<0.\quad (1)$\\
		Do $m<-5$ nên $f(-1)=-m-5>0.\quad (2)$\\ 
		Mặt khác  $f(0)=m-3<0\quad (3)$\\
		$\lim\limits_{x\to  +\infty}f(x)=+\infty$ nên tồn tại $b>0$ sao cho $f(b)>0.\quad (4)$\\
		Từ $(1)$ và $(2)$, suy ra phương trình có nghiệm thuộc khoảng  $(-\infty;-1)$.\\
		Từ $(2)$ và $(3)$, suy ra phương trình có  nghiệm thuộc khoảng $(-1;0)$.\\
		Từ $(3)$ và $(4)$, suy ra phương trình có  nghiệm thuộc khoảng $(0;+\infty)$.\\
		Vậy $m<-5$ thỏa mãn yêu cầu  bài toán.\\
		 Mà $\heva{&-10<m<10\\&m\in \mathbb{Z}}\Rightarrow m\in \left\lbrace -9;-8;-7;-6 \right\rbrace $.
	} 
\end{ex}

% nguồn: Ok_TracNghiemLop11_Lan1
\begin{ex}%[1D3K3-5]%[ma:e95177b1]
	 Tính tổng $T=15+20+25+\cdots+7515$.
	 \choice
		 { $T=5651625$}
		{ \True $T=5651265$} 
		{ $T=5651526$} 
		{  $T=5651256$}
	\loigiai{ 
		Ta thấy các số hạng của tổng $T$ tạo thành một cấp số cộng với  số hạng đầu ${{u}_{1}}=15$ và công sai $d=5$. Giả sử tổng trên có $n$ số hạng thì ${{u}_{n}}=7515$\\
		$\Leftrightarrow {{u}_{1}}+(n-1)d=7515\Leftrightarrow  15+(n-1)5=7515\Leftrightarrow n=1501$.\\
	Vậy $T={{S}_{1501}}=\dfrac{(2{{u}_{1}}+1500d)\cdot 1501}{2}=\dfrac{(2\cdot  15+1500\cdot 5)\cdot 1501}{2}=5651265$.}
\end{ex}

% nguồn: Ok_TracNghiemLop11_Lan1
\begin{ex}%[1D3K3-5]%[ma:21000d03]
	Xét các số nguyên dương chia hết cho $3$. Tổng số $50$ số nguyên dương đầu tiên đó bằng: 
	 \choice
		{ \True $3825$} 
		 { $3900$}
		{ $7500$} 
		{ $7650$} 
	 \loigiai{
		Số  nguyên dương chia hết cho $3$ có dạng $3n\,\,(n\in {{\mathbb{N}}^{*}})$ nên chúng lập thành cấp số cộng $${{u}_{n}}=3n\Rightarrow \heva{
		&  {{u}_{1}}=3 \\
		& {{u}_{50}}=150 \\ 
		}\Rightarrow  {{S}_{50}}=\dfrac{50}{2}({{u}_{1}}+{{u}_{50}})=3825.$$
	\textbf{Chú ý:} ${{S}_{n}}=\dfrac{n}{2}({{u}_{1}}+{{u}_{n}})=n{{u}_{1}}+\dfrac{n(n-1)}{2}d$.} 
\end{ex}

% nguồn: Ok_TracNghiemLop11_Lan1
\begin{ex}%[1D3K3-5]%[ma:9785bcca]
	Tổng $n$ số hạng đầu tiên của một cấp số cộng  là ${{S}_{n}}=\dfrac{3{{n}^2}-19n}{4}$ với $n\in {{\mathbb{N}}^{*}}$. Tìm số hạng đầu tiên ${{u}_{1}}$ và công sai $d$ của cấp số cộng đã cho.
	 \choice
		 { \True ${{u}_{1}}=-4;d=\dfrac{3}{2}$}
		{ ${{u}_{1}}=\dfrac{5}{2};d=\dfrac{1}{2}$} 
		{ ${{u}_{1}}=-\dfrac{3}{2};d=-2$} 
		{  ${{u}_{1}}=2;d=-\dfrac{1}{2}$}
	 \loigiai{
		Ta  có $$\dfrac{3{{n}^2}-19n}{4}=\dfrac{3}{4}{{n}^2}-\dfrac{19}{4}n={{S}_{n}}=n{{u}_{1}}+\dfrac{{{n}^2}-n}{2}d=\dfrac{d}{2}{{n}^2}+\left({{u}_{1}}-\dfrac{d}{2}\right)n
		\Leftrightarrow \heva{ 
		& \dfrac{d}{2}=\dfrac{3}{4} \\ 
		 & {{u}_{1}}-\dfrac{d}{2}=-\dfrac{19}{4} \\
		}\Leftrightarrow  \heva{
		 & {{u}_{1}}=-4 \\
		&  d=\dfrac{3}{2}. \\
	 }$$}
\end{ex}

% nguồn: Ok_TracNghiemLop11_Lan1
\begin{ex}%[1D3K3-5]%[ma:2b58b276]
	Tổng $n$ số hạng đầu tiên của một cấp  số cộng là ${{S}_{n}}={{n}^2}+4n$ với $n\in {{\mathbb{N}}^{*}}$. Tìm số hạng tổng quát ${{u}_{n}}$ của cấp số cộng đã cho.
	 \choice
		 { ${{u}_{n}}=5\cdot {{\left(\dfrac{8}{5} \right)}^{n-1}}$}
		{ ${{u}_{n}}=3n+2$} 
		{ ${{u}_{n}}={{5\cdot  3}^{n-1}}$}
		{  \True ${{u}_{n}}=2n+3$}
	 \loigiai{
		Ta có ${{n}^2}+4n={{S}_{n}}=\dfrac{d}{2}{{n}^2}+\left({{u}_{1}}-\dfrac{d}{2}\right)n\Leftrightarrow  \heva{
		& \dfrac{d}{2}=1 \\ 
		& {{u}_{1}}-\dfrac{d}{2}=4 \\ 
		}$ 
		$\Leftrightarrow \heva{ 
		 & {{u}_{1}}=5 \\
		 & d=2
	}\Rightarrow {{u}_{n}}=2n+3$.} 
\end{ex}


%%%%%%%%%%%%%%%% ĐÚNG/SAI %%%%%%%%%%%%%%%%

% nguồn: Ok_dungsai_lop11_Lan1
\begin{ex}%[1D3H3-5]%[ma:97aadecf]
 Cho hàm số $y=f(x)$ liên tục trên đoạn $[1;5]$ và $f(1)=2$, $f(5)=10$. Xét tính đúng sai của  các khẳng định sau?
 \choiceTF 
  {Phương trình  $f(x)=6$ vô nghiệm}
  {\True Phương trình $f(x)=7$ có ít  nhất một nghiệm trên khoảng $(1;5)$}
  {Phương trình $f(x)=2$ có hai  nghiệm $x=1$, $x=5$}
  {Phương trình $f(x)=7$  vô nghiệm}
  \loigiai{
 Vì hàm số $y=f(x)$ liên tục trên đoạn $[1;5]$ và $f(1)=2$, $f(5)=10$ nên phương trình $f(x)=a$ với $a \in (2;  10)$ luôn có ít nhất một nghiệm thuộc khoảng $(1;5)$.
 \begin{itemchoice}
  \itemch Vì $6 \in (2; 10)$  nên phương trình $f(x)=6$ có nghiệm.
  \itemch Vì $7\in (2;10)$ nên phương trình $f(x)=7$ có ít nhất một nghiệm  trên khoảng $(1;5)$.
  \itemch Do  $f(1)=2\neq 0$, $f(5)=10 \neq 0$ nên $x=1$, $x=5$ không là nghiệm của phương trình.
  \itemch  Phương trình $f(x)=7$ có nghiệm do $7\in (2;10)$.
 \end{itemchoice}
  }
\end{ex}

% nguồn: Ok_dungsai_lop11_Lan1
\begin{ex}%[1D3H3-5]%[ma:c417dc69]
 Xét tính đúng sai của các khẳng định  sau?
 \choiceTF 
  {\True Phương trình $x^4+mx^2-2mx-2=0$ luôn có nghiệm với mọi giá trị của tham số $m$} 
  {Phương trình $3x^6-3x^3+5x-2=0$ \textbf{không} có nghiệm  thuộc khoảng $(-2;2)$}
  {\True Phương trình $x^3-3x+1=0$ có ba nghiệm  phân biệt}
  {\True Phương  trình $m(x-1)^2(x-2)+2x-3=0$ luôn có nghiệm với mọi giá trị của tham số $m$}
 \loigiai{ 
  \begin{itemchoice}
   \itemch Hàm số $f(x)=x^4+mx^2-2mx-2$ liên tục trên đoạn $[0;2]$, có  $f(0)\cdot f(2) = -2\cdot 14 = -28<0$ nên phương trình $f(x)=0$ có ít nhất một nghiệm thuộc khoảng $(0;2)$ với mọi $m$ hay phương trình $x^4+mx^2-2mx-2=0$ luôn có nghiệm với mọi giá trị của tham số $m$.
   \itemch Hàm số $g(x)=3x^6-3x^3+5x-2$ liên tục  trên đoạn $[0;1]$, có $g(0)\cdot g(1) = -2\cdot 3 =-6<0$ nên phương trình $g(x)=0$ có ít nhất một nghiệm thuộc khoảng $(0;1)$ hay phương trình $g(x)=0$ có nghiệm thuộc khoảng $(-2;2)$.
   \itemch Hàm số $h(x)=x^3-3x+1$ liên tục trên các đoạn $[-2;0]$, $[0;1]$, $[1;2]$ và có  $h(-2)=-1$, $h(0)=1$, $h(1)=-1$, $h(2)=3$.\\
   Do $h(-2)\cdot h(0)<0$, $h(0)\cdot h(1)<0$, $h(1)\cdot h(2)<0$ nên phương trình $h(x)=0$ có ba nghiệm $x_1\in (-2;0)$, $x_2\in (0;1)$, $x_3\in (1;2)$ là  ba nghiệm phân biệt.\\
   Mặt khác phương trình $h(x)=0$ là phương trình bậc  ba nên có không quá ba nghiệm.\\
   Vậy phương trình $h(x)=0$ có ba nghiệm  phân biệt.
   \itemch Hàm số $p(x)=m(x-1)^2(x-2)+2x-3$ liên tục  trên đoạn $[1;2]$, có $p(1)\cdot p(2) = (-1)\cdot 1 = -1<0$ nên phương trình $p(x)=0$ luôn có nghiệm thuộc khoảng $(1;2)$ với mọi $m$ hay phương trình $m(x-1)^2(x-2)+2x-3=0$ luôn có nghiệm với mọi giá trị của tham số $m$.
  \end{itemchoice}
 } 
\end{ex}

% nguồn: Ok_dungsai_lop11_Lan1
\begin{ex}%[1D3V3-5]%[ma:c113d51d]
 Cho phương trình $x^4+8x^3+11x^2-32x-60=0$. Xác định  tính đúng sai của các khẳng định sau.
 \choiceTF 
  {Phương trình đã cho chỉ có một nghiệm trong  khoảng $(-5; 5)$}
  {\True Phương trình đã cho có  hai nghiệm trong khoảng $(-3; 3)$}
  {\True Phương trình đã cho có một nghiệm trong khoảng $(-3; 0)$} 
  { Phương trình đã cho  không có nghiệm trong khoảng $(0; 3)$}
 \loigiai{ 
  Đặt $f(x)=x^4+8x^3+11x^2-32x-60$ là hàm số liên tục trên $\mathbb{R}$. Ta có  $f(-5)=0$, $f(-3)=0$, $f(-2)=0$, $f(0)=-60$, $f(3)=240$.
   \\
  Do đó $f(x)=0$ có 3 nghiệm là $-5, -3, -2$ và có một nghiệm thuộc khoảng $(0;3)$. 
  \begin{itemchoice}
   \itemch  Sai vì trong khoảng $(-5; 5)$ phương trình đã cho có $3$ nghiệm.
   \itemch Đúng vì  trong khoảng $(-3;3)$ phương trình đã cho có hai nghiệm
   \itemch Đúng vì trong khoảng $(-3;0)$ phương trình đã cho có  một nghiệm.
   \itemch Sai vì trong khoảng  $\left(0;3\right)$ có một nghiệm.
  \end{itemchoice}
 } 
\end{ex}


%%%%%%%%%%%%%%%% TRẢ LỜI NGẮN %%%%%%%%%%%%%%%%

% nguồn: Ok_TraLoiNgan_Lop11_Lan2
\begin{ex}%[1D3V3-5]%[ma:afbb97b6]
	Cho các số thực $a$, $b$, $c$ thỏa  mãn $\heva{
	& a-b+c-1>0\\ 
	&  4a+2b+c+8<0
	}$.  Tìm số nghiệm của phương trình $x^3+ax^2+bx+c=0$.
	\shortans{3} 
	 \loigiai{
		Xét hàm số  $f(x)=x^3+ax^2+bx+c$, là hàm đa thức nên liên tục trên $\mathbb{R}$.
		\begin{itemize}
			\item Ta có $\lim\limits_{x\to-\infty}f(x)=-\infty$ và $f(-1)=a-b+c-1>0$ nên phương trình $f(x)=0$ có ít  nhất một nghiệm trong khoảng $(-\infty;-1)$.
			\item $f(-1)\cdot f(2)=(a-b+c-1)(4a+2b+c+8)<0$ nên phương trình  $f(x)=0$ có ít nhất một nghiệm trong khoảng $(-1;2)$.
			\item $f(2)=4a+2b+c+8<0$ và  $\lim\limits_{x\to+\infty}f(x)=+\infty$ nên phương trình $f(x)=0$ có ít nhất một nghiệm trong khoảng $(2;+\infty)$.
		\end{itemize}
		Mà phương trình là bậc $3$ nên có tối đa $3$  nghiệm.\\
		Vậy phương trình có  $3$ nghiệm.
	 }
\end{ex}

% nguồn: Ok_TraLoiNgan_Lop11_Lan2
\begin{ex}%[1D3V3-5]%[ma:24dba841]
	Cho các số thực $a$, $b$, $c$ thỏa mãn $4a+c > 8+2b$ và $a+b+c <-1$. Khi đó số nghiệm thực  phân biệt của phương trình $x^3+ax^2+bx+c=0$ bằng
	\shortans{3} 
	 \loigiai{Đặt $f(x)=x^3+ax^2+bx+c$. \\
		Tập xác định $\mathscr{D}=\mathbb{R}$.\\ 
		Từ giả thiết $\heva{&4a+c > 8+2b\Rightarrow-8+4a-2b+c > 0\Rightarrow f\left(-2\right) > 0  \\
		&a+b+c <-1\Rightarrow 1+a+b+c <  0\Rightarrow f(1) < 0..}$\\
		 Do đó $f\left(-2\right)\cdot f(1) < 0$ nên phương trình có ít nhất một nghiệm trong $\left(-2;1\right)$.\\
		Ta  nhận thấy
		$\lim\limits_{x\to-\infty} f(x)=-\infty$ mà $f\left(-2\right) > 0$ nên phương trình có ít nhất một nghiệm $\alpha \in  \left(-\infty;-2\right)$.\\
		Tương tự $\lim\limits_{x\to+\infty} f(x)=+\infty$ mà $f(1) < 0$ nên phương trình có ít nhất một nghiệm  $\beta \in \left(1;+\infty \right)$.\\
		Vì $f\left(x\right)$ là đa thức bậc ba có tối đa ba nghiệm thực phân biệt nên số nghiệm của phương trình đã cho có đúng  $3$ nghiệm thực phân biệt.
	} 
\end{ex}

% nguồn: Ok_TraLoiNgan_Lop11_Lan2
\begin{ex}%[1D3V3-5]%[ma:92a4c799]
	Cho phương trình $4x^4 + 2x^2 - x - 3 = 0$.  Phương trình có đúng $m$ nghiệm trên khoảng $(-1;1)$. Tìm $m$.
	\shortans{2} 
	 \loigiai{
		Xét $f(x)=4x^4+2x^2-x-3$ trên đoạn $[-1;1]$. Ta có $f(x)$ liên tục trên đoạn  này.\\
		$f(-1)=4,\; f(0)=-3,\; f(1)=2 \Rightarrow  f(-1)f(0)<0,\; f(1)f(0)<0$.\\
		Do  đó phương trình $f(x)=0$ có hai nghiệm trong khoảng $(-1;1)$.\\
		Ta có $f'(x)=6x^3+4x-1$. Khi đó $f'(-1)=-11,\; f'(1)=9 \Rightarrow  f'(-1)f'(1)<0$.\\
		Do đó phương trình  $f'(x)=0$ có nghiệm trong khoảng $(-1;1)$.\\
		Vì $f''(x)=18x^2+4>0$ với mọi $x\in(-1;1)$  nên $f'(x)$ đồng biến trên khoảng này $\Rightarrow$ phương trình $f'(x)=0$ có duy nhất một nghiệm trong khoảng $(-1;1)$.\\
		Suy ra  $f(x)=0$ có tối đa hai nghiệm trên khoảng $(-1;1)$.\\
		Vậy phương trình  có đúng hai nghiệm trên khoảng $(-1;1)$.
	 }
\end{ex}

% nguồn: Ok_TraLoiNgan_Lop11_Lan2
\begin{ex}%[1D3V3-5]%[ma:96541364]
	Cho số thực  $a$, $b$, $c$ thỏa mãn $\heva{&{-8+4a-2b+c > 0} \\
	 &{8+4a+2b+c < 0}.
	}$ Số giao điểm của đồ thị hàm số $f\left(x\right)=x^3+ax^2+bx+c$  với trục hoành là
	 \shortans{3}
	 \loigiai{Đặt $f(x)=x^3+ax^2+bx+c$. Ta có $f(x)$ là hàm đa thức nên liên tục trên $\mathbb{R}$.\\
		Khi  đó $\heva{&{f(2)=8+4a+2b+c < 0} \\
		&{f\left(-2\right)=-8+4a-2b+c  > 0}
		}\Rightarrow  f\left(-2\right).f(2) < 0$ $\Rightarrow$ đồ thị hàm số $y=f(x)$ cắt trục $Ox$ tại ít nhất một điểm trong khoảng $\left(-2; 2\right)$.\\
		Mặt  khác, ta có
		$\heva{&{f(2) < 0}  \\
		&{\lim\limits_{x\to+\infty} f(x)=+\infty} 
		}$ $\Rightarrow$ đồ thị hàm số $y=f(x)$ cắt trục $Ox$ tại ít  nhất một điểm trong khoảng $\left(2;+\infty \right)$.\\
		Lại  có
		$\heva{&{f\left(-2\right) > 0}  \\
		&{\lim\limits_{x\to-\infty} f(x)=-\infty} 
		}$ $\Rightarrow$ đồ  thị hàm số $y=f(x)$ cắt trục $Ox$ tại ít nhất một điểm trong khoảng $\left(-\infty;-2\right)$.\\
		 Vì hàm số $f(x)$ là hàm bậc ba nên đồ thị của nó cắt trục $Ox$ tối đa tại ba điểm.\\
		Vậy đồ thị hàm số $y=f(x)$ cắt trục $Ox$ tại đúng $3$  điểm.
	 }
\end{ex}


%%%%%%%%%%%%%%%% TỰ LUẬN %%%%%%%%%%%%%%%%

% nguồn: Ok_TuanLuan_Toan11_Lan2
\begin{ex}%[1D3H3-5]%[ma:fe90d6e4]
	Chứng minh rằng phương trình sau đây có nghiệm: $x^4-3x+1=0$ 
	 \loigiai{
		Xét  hàm số $f(x)=x^4-3x+1$ liên tục trên $\mathbb{R}$ nên cũng liên tục trên đoạn $\left[0;1 \right]$.\\
		 Ta có $f(0)=1 > 0$, $f(1)=-1 < 0 \Rightarrow f(0)\cdot f(1) < 0$.\\
		Vậy phương trình có  ít nhất một nghiệm trên khoảng $(0;1)$.
	 }
\end{ex}

% nguồn: Ok_TuanLuan_Toan11_Lan2
\begin{ex}%[1D3H3-5]%[ma:edb95d0c]
	Chứng minh rằng phương trình $x^3+(m+3)x^2+(1-m)x-1=0$ luôn có nghiệm với mọi  giá trị của $m$.
	\loigiai{ 
		Xét hàm  số $f(x)=x^3+(m+3)x^2+(1-m)x-1$ liên tục trên $\mathbb{R}$ nên cũng liên tục trên khoảng $(0;1)$.\\
		Ta có 
		$\begin{aligned}
		& f(0)=-1 <  0 \\
		& f(1)=4 >  0 .
	\end{aligned}$\\
	Suy ra  $f(0)\cdot f(1)<0$.\\
	Vậy trên  khoảng $(0;1)$ phương trình có ít nhất một nghiệm.
} 
\end{ex}

% nguồn: Ok_TuanLuan_Toan11_Lan2
\begin{ex}%[1D3H3-5]%[ma:9635168a]
	Chứng minh rằng phương  trình $2x^3-6x+1=0$ có 3 nghiệm trong khoảng $(-2;2)$.
	\loigiai{ 
		Xét hàm số $f(x)=2x^3-6x+1$ liên tục trên $\mathbb{R}$ nên cũng liên tục trên các đoạn $[-2;-1]$; $[-1;1]$;  $[1;2]$.\\
		 Ta có
		$\begin{aligned}
		& f(-2)=-3 < 0  \\
		& f(-1)=5 > 0 \\ 
		& f(1)=-3 \\ 
		& f(2)=5  .
	\end{aligned}$\\
	Vậy trên mỗi khoảng $(-2;-1)$; $(-1;1)$; $(1;2)$ phương trình có ít nhất nghiệm. Mà ba khoảng này rời nhau nên phương trình có  $3$ nghiệm trong khoảng $(-2;2)$.
} 
\end{ex}

% nguồn: Ok_TuanLuan_Toan11_Lan2
\begin{ex}%[1D3H3-5]%[ma:94fd1d61]
	Tìm tất cả các giá  trị của tham số $m$ để phương trình $(m^2-3m+2)x^3-3x+1=0$ có nghiệm.
	\loigiai{ 
		Đặt $\left\{\begin{aligned}
		& f(x)=(m^2-3m+2)x^3-3x+1 \\ 
		& A=m^2-3m+2  \\
	\end{aligned} \right.$ \\
	Xét $A=m^2-3m+2=0 \Leftrightarrow $  $m=1$ hoặc $m=2$.\\
	Khi đó phương  trình trở thành $-3x+1=0 \Leftrightarrow x=\dfrac{1}{3}$. \\
	Xét $A=m^2-3m+2\ne  0 \Leftrightarrow $ $m\ne 1$ và $m\ne 2$.\\
	Khi đó: 
	\begin{itemize}
		\item Xét hàm số  $f(x)=(m^2-3m+2)x^3-3x+1$ đây là hàm đa thức, xác định trên $\mathbb{R}$ nên liên tục trên $\mathbb{R}$.
		 \item Mặt khác, ta có:
		\begin{itemize}[]
			\item [\bf TH 1.]  $A>0 \Rightarrow m\in (-\infty;1)\cup (2;+\infty)$. \\
			$\lim\limits_{x\to+\infty}f(x) =\lim\limits_{x\to+\infty}\left[Ax^3-3x+1 \right] =+\infty$ nên  tồn tại $x_1\in \mathbb{R}$ sao cho $f(x_1)>0$.\\
			$\lim\limits_{x\to-\infty}f(x) =\lim\limits_{x\to-\infty}\left[Ax^3-3x+1 \right] =-\infty$ nên tồn tại $x_2\in \mathbb{R}$ sao cho $f(x_2) <  0$.\\
			Áp dụng hệ  quả của định lí về giá trị trung gian, suy ra tồn tại $t\in (x_2;x_1)$ sao cho $f(t)=0$.\\
			\item [\bf  TH 2.] $A< 0\Rightarrow m\in (1;2)$. \\
			$\lim\limits_{x\to+\infty}f(x)  =\lim\limits_{x\to+\infty}\left[Ax^3-3x+1 \right]=-\infty$ nên tồn tại $x_1\in \mathbb{R}$ sao cho $f(x_1)<0$.\\
			$\lim\limits_{x\to-\infty}f(x) =\lim\limits_{x\to-\infty}\left[Ax^3-3x+1 \right]=+\infty$ nên tồn tại $x_2\in \mathbb{R}$  sao cho $f(x_2) > 0$.\\
			Áp dụng hệ quả của định lí về giá trị trung gian, suy ra tồn  tại $t\in (x_2;x_1)$ sao cho $f(t)=0$.
		\end{itemize}
		Vậy  phương trình đã cho có nghiệm với mọi $m\in \mathbb{R}$.
	\end{itemize}
	 }
\end{ex}

% nguồn: Ok_TuanLuan_Toan11_Lan2
\begin{ex}%[1D3H3-5]%[ma:6aaf56e4]
	Với mọi giá trị thực của tham số $m$ chứng minh phương trình \begin{align*}
	 (m^2+2m+3)(x^4-5x^2+4)+4x^3-9x=0
\end{align*} luôn có ít nhất ba nghiệm thực.
 \loigiai{
	Đặt $f(x)=(m^2+2m+3)(x^4-5x^2+4)+4x^3-9x$ là hàm số  liên tục trên $\mathbb{R}$.\\
	Các  giá trị
	$\begin{aligned}
	& f(-2)=-14 \\ 
	& f(-1)=5  \\
	 & f(1)=-5 \\
	& f(2)=14  .
	\end{aligned}$\\
	Ta có $f(-2)\cdot f(-1)<0$, $f(-1)\cdot  f(1)<0$, $f(1)\cdot f(2)<0$ với mọi $m\in \mathbb{R}$
	$\Rightarrow f(x)=0$ luôn có ít nhất $1$  nghiệm trên mỗi khoảng $(-2;-1)$, $(-1;1)$, $(1;2)$.
 }
\end{ex}

% nguồn: Ok_TuanLuan_Toan11_Lan2
\begin{ex}%[1D3H3-5]%[ma:1e5d8309]
	Tìm tất cả các giá trị của tham số thực sao cho phương trình sau có  nghiệm:\break $(2m^2-5m+2)(x-1)^{2017}(x^{2018}-2)+2x+3=0$.
	 \loigiai{
		\begin{itemize}
			\item Nếu $2m^2-5m+2=0$ thì phương trình đã cho trở thành $2x+3=0  \Leftrightarrow x=-\dfrac{3}{2}$.
			\item Nếu $2m^2-5m+2\ne 0$ phương trình đã cho là một đa thưc bậc  lẻ nên phương trình có ít nhất một nghiệm.
		\end{itemize}
		 Vậy với mọi phương trình đã cho luôn có ít nhất một nghiệm.
	} 
\end{ex}

% nguồn: Ok_TuanLuan_Toan11_Lan2
\begin{ex}%[1D3H3-5]%[ma:18da3aab]
	Chứng minh rằng với mọi  giá trị thực của tham số $m$ phương trình sau luôn có nghiệm
	 $x^3-2x^2+x-1-\dfrac{2}{x}+\dfrac{m}{x+3}=0.$
	\loigiai{ 
		Tập  xác định $\mathscr{D}=R\backslash \{-3;0\}$.\\
		Trên $\mathscr{D}$, phương trình đã cho tương đương  với phương trình:
		 $$mx+(x-2)(x+3)(x^3+x+1)=0.\,(*)$$
		Đặt $f(x)=mx+(x-2)(x+3)(x^3+x+1)$.\\ 
		Nhận xét $f$ liên tục trên $\mathbb{R}$ và $\left\{\begin{aligned}
		& f(-3)\cdot f(0)=18m  \\ & f(0)\cdot f(2)=-12m. \\
	\end{aligned} \right.$\\
	\textbf{TH1:}  Nếu $m=0$, phương trình luôn có $1$ nghiệm $x=2$.\\
	\textbf{TH2:} Nếu $m\ne  0 \Rightarrow (*)$ có ít nhất một nghiệm thuộc $(-3;0)$.\\
	Suy ra $\forall m\in \mathbb{R}$, luôn có ít nhất một nghiệm thuộc $\mathscr{D}$.\\ 
	Vậy $\forall m\in \mathbb{R}$, phương trình đã cho luôn có ít nhất  một nghiệm.
	} 
\end{ex}

% nguồn: Ok_TuanLuan_Toan11_Lan2
\begin{ex}%[1D3H3-5]%[ma:34b122e7]
	Vậy với mọi số thực $m$ thì phương trình đã cho có ít nhất $3$ nghiệm thực. Chứng  minh rằng phương trình $(1-m^2)x^5-3x-1=0$ luôn có nghiệm với mọi giá trị của tham số $m$.
	 \loigiai{
		 Xét hàm số $f(x)=(1-m^2)x^5-3x-1$.\\
		Vì $f(0)=-1 < 0$ và $f(-1)=m^2+1 >  0$ nên $f(0)\cdot f(-1)<0$ với mọi $m$.\\
		Mặt khác $f(x)$ là hàm đa  thức, liên tục trên $\mathbb{R}$ nên liên tục trên đoạn $[-1;0]$.\\
		Từ và suy ra phương  trình $f(x)=0$ có ít nhất một nghiệm trong khoảng $(-1;0)$ nghĩa là phương trình $(1-m^2)x^5-3x-1=0$ luôn có nghiệm với mọi $m$.
	 }
\end{ex}

% nguồn: Ok_TuanLuan_Toan11_Lan2
\begin{ex}%[1D3H3-5]%[ma:418e6380]
	Tìm tất cả các giá  trị của tham số $m$ sao cho phương trình $x^3-3x^2+(2m-2)x+m-3=0$ có ba nghiệm $x_1$, $x_2$, $x_3$ thỏa mãn $x_1 <-1 < x_2 < x_3$.
	 \loigiai{
		Đặt $f(x)=x^3-3x^2+(2m-2)x+m-3$.\\ 
		Ta thấy hàm số liên  tục trên $\mathbb{R}$.\\
		 \textbf{Điều kiện cần:} $af(-1) > 0 \Leftrightarrow-m-5 > 0 \Leftrightarrow m <-5$.\\
		\textbf{Điều kiện đủ:}  với $m <-5$ ta có\\
		\begin{itemize}
			\item $\lim\limits_{x\to-\infty}f(x)=-\infty$ nên tồn tại $a <-1$  sao cho $f(a) < 0$. \\
		\end{itemize}
		Mặt khác $f(-1)=-m-5 > 0$. Suy ra $f(a)\cdot f(-1) < 0$.\\ 
		Do đó tồn tại $x_1\in  (a;-1)$ sao cho $f(x_1)=0$.\\
		\begin{itemize}
			\item $f(0)=m-3 < 0$, $f(-1) > 0$. Suy ra $f(0)\cdot f(-1)  < 0$.\\
		\end{itemize}
		Do đó tồn tại $x_2\in (-1;0)$ sao  cho $f(x_2)=0$.\\
		\begin{itemize}
			\item $\lim\limits_{x\to+\infty}f(x)=+\infty$ nên tồn tại $b >  0$ sao cho $f(b) > 0$. \\
		\end{itemize}
		Mặt khác $f(0)<0$. Suy ra  $f(0)\cdot f(b)<0$.\\
		Do đó tồn tại $x_3\in  (0;b)$ sao cho $f(x_3)=0$.\\
		Vậy $m <-5$ thỏa mãn yêu  cầu bài toán.
	 }
\end{ex}

% nguồn: Ok_TuanLuan_Toan11_Lan2
\begin{ex}%[1D3H3-5]%[ma:1c1daeae]
	Cho phương trình $(m^2-m+5)x^3-3x-6=0$. Chứng minh rằng: Với mọi $m\in \mathbb{R}$, phương  trình đã cho có ít nhất $1$ nghiệm.
	\loigiai{ 
		Xét hàm số $f(x)=(m^2-m+5)x^3-3x-6$ trên  $\mathbb{R}$.\\
		 Ta có
		$\begin{aligned}
		& f(0)=-6 <  0\\
		& f(2)=8m^2-8m+28= 2(2m-1)^2+26 >  0, \forall m\in \mathbb{R}.
	\end{aligned}$\\
	Hàm số $y=f(x)$ liên tục trên đoạn $[0;2]$ và  $f(0)\cdot f(2) < 0,\forall m\in \mathbb{R} $.\\
	Suy ra phương trình đã cho có ít nhất một nghiệm trong khoảng  $(0;2) $.
} 
\end{ex}

% nguồn: Ok_TuanLuan_Toan11_Lan2
\begin{ex}%[1D3H3-5]%[ma:59ad60ca]
	Với mọi giá trị thực của  tham số $m$ chứng minh phương trình $(1-m^2)x^5-3x-1=0$ luôn có nghiệm thực.
	\loigiai{ 
		Xét hàm số $f(x)=(1-m^2)x^5-3x-1$ là hàm đa thức nên liên tục  trên $\mathbb{R}$.\\
		 Các giá trị
		$\begin{aligned}
		&  f(0)=-1 \\
		& f(-1)=m^2+1 . 
	\end{aligned}$\\
	Ta có $f(0)\cdot f(-1)=-1\cdot (m^2+1)<0,\forall m$ nên phương trình  $f(x)=0$ có ít nhất một nghiệm thực thuộc $(-1;0)$.\\
	Vậy với mọi  giá trị thực của tham số $m$ phương trình $(1-m^2)x^5-3x-1=0$ luôn có nghiệm thực.
} 
\end{ex}

% nguồn: Ok_TuanLuan_Toan11_Lan2
\begin{ex}%[1D3H3-5]%[ma:44867de1]
	Cho phương trình $ax^2+bx+505c=0$ ($a\ne 0$) thỏa mãn  $a+2b+2022c=0$. Chứng minh phương trình trên có nghiệm.
	\loigiai{ 
		 Đặt $f(x)=ax^2+bx+505c$ với $a\ne 0$. \\
		Hàm số  $f(x)$ có tập xác định $\mathbb{R}$.\\
		Vì $f(x)$ là hàm số đa thức bậc hai liên tục trên $\mathbb{R}$ nên liên tục trên đoạn $\left[0;\dfrac{1}{2}  \right]$. \tagEX{1}
		\noindent  Lại có: $f(0)=505c$, $f\left(\dfrac{1}{2} \right) =\dfrac{1}{4}a+\dfrac{1}{2}b+505c =\dfrac{1}{4}(a+2b+2022c)-\dfrac{1}{2}c =-\dfrac{1}{2}c$.\\
		Khi  đó: $f(0)\cdot f\left(\dfrac{1}{2} \right) =-\dfrac{505}{2}c^2\le 0$.
		\begin{itemize}[\textbf{TH.1:}]
			\item Nếu  $c=0$ thì phương trình $f(x)=0$ có hai nghiệm phân biệt $x=0$; $x=-\dfrac{b}{a}$.
			\item Nếu $c\ne 0$ thì $f(0)\cdot  f\left(\dfrac{1}{2} \right) < 0$. \tagEX{2}
		\end{itemize}
		Từ $(1)$ và $(2)$  suy ra phương trình $f(x)=0$ có ít nhất một nghiệm thuộc khoảng $\left(0;\dfrac{1}{2} \right)$.\\
		Vậy phương  trình đã cho luôn có nghiệm.
	} 
\end{ex}

% nguồn: Ok_TuanLuan_Toan11_Lan2
\begin{ex}%[1D3V3-5]%[ma:f920ef3e]
	Chứng minh rằng với  mọi giá trị của tham số $m$ phương trình \break$(1-m^2)(x+1)^3+x^2-x-3=0$ luôn có nghiệm.
	\loigiai{ 
		Xét hàm số $f(x)=(1-m^2)(x+1)^3+x^2-x-3$ là  hàm số liên tục trên $\mathbb{R}$ nên nó liên tục trên $[-2;-1]$.\\
		Ta có $f(-1)=-1; f(-2)=m^2+2$ nên $f(-1) \cdot  f(-2) < 0, \forall m \in \mathbb{R}$.\\
		Khi  đó phương trình $f(x)=0$ có ít nhất một nghiệm trong khoảng $(-2;-1)$.\\
		Vậy phương trình $(1-m^2)(x+1)^3+x^2-x-3=0$ luôn có ít nhất một  nghiệm với mọi giá trị của $m$.
	 }
\end{ex}

% nguồn: Ok_TuanLuan_Toan11_Lan2
\begin{ex}%[1D3V3-5]%[ma:e2ed179b]
	Chứng minh rằng với mọi  giá trị của tham số $m$ phương trình $(1-m^2)x^5-3x-1=0$ luôn có nghiệm.
	 \loigiai{
		Xét hàm số $f(x)=(1-m^2)x^5-3x-1$ là hàm số liên tục trên $\mathbb{R}$ nên nó liên tục trên  $[-1; 0]$.\\
		Ta có $f(0)=-1$ và $f(-1)=m^2+1$ nên $f(0) \cdot f(-1) < 0, \forall m \in  \mathbb{R}$.\\
		Khi đó phương trình  $f(x)=0$ có ít nhất một nghiệm trong khoảng $(-1; 0)$.\\
		Vậy phương trình $(1-m^2)x^5-3x-1=0$ luôn có ít nhất một  nghiệm.
	} 
\end{ex}

% nguồn: Ok_TuanLuan_Toan11_Lan2
\begin{ex}%[1D3V3-5]%[ma:a0309f15]
	Chứng minh rằng phương trình $\dfrac{1}{\cos x}-\dfrac{1}{\sin x}=m$ luôn có nghiệm  với mọi giá trị của tham số $m$.
	\loigiai{ 
		 Điều kiện: $x \neq k\dfrac{\pi}{2}, k \in \mathbb{Z}$.\\
		Xét hàm số $f(x)=\sin x-\cos x-m\sin  x \cos x$ liên tục trên $\left[0; \dfrac{\pi}{2}\right]$ và $f(0)\cdot f\left(\dfrac{\pi}{2}\right)=-1< 0$.\\
		Do đó phương trình $f(x)=0$ có ít  nhất một nghiệm $x_0\in \left(0; \dfrac{\pi}{2}\right) \Rightarrow x_0\neq k\dfrac{\pi}{2}$.\\
		Vậy phương trình  đã cho có ít nhất một nghiệm.
	} 
\end{ex}

% nguồn: Ok_TuanLuan_Toan11_Lan2
\begin{ex}%[1D3V3-5]%[ma:3513c087]
	Chứng minh rằng phương trình $\cos x = x$ có nghiệm. 
	\loigiai{ 
		Ta có $\cos x = x \Leftrightarrow  \cos x - x = 0$.\\
		Xét hàm số $f(x) = \cos x - x$ là hàm số liên tục  trên $\mathbb{R}$ nên nó liên tục trên $\left[0; \dfrac{\pi}{2}\right]$.\\
		Ta có $f(0)  = 1$ và $f\left(\dfrac{\pi}{2}\right) = -\dfrac{\pi}{2}$ nên $f(0) \cdot f\left(\dfrac{\pi}{2}\right) < 0$.\\
		Khi đó phương  trình $f(x) = 0$ có ít nhất một nghiệm trong khoảng $\left(0; \dfrac{\pi}{2}\right)$.
	} 
\end{ex}

% nguồn: Ok_TuanLuan_Toan11_Lan2
\begin{ex}%[1D3V3-5]%[ma:97e5eae1]
	Chứng minh rằng phương trình $a\cos 3x+b\cos 2x+c\cos x+\sin x=0$ luôn có nghiệm trên  $[0; 2\pi]$.
	\loigiai{ 
		Xét hàm  số $f(x)=a\cos 3x+b\cos 2x+c\cos x+\sin x$ liên tục trên $\mathbb{R}$.\\
		Nên cũng liên tục trên các đoạn  $[0; 2\pi]$.\\
		$f(0)=a+b+c$, $f\left(\dfrac{\pi}{2}\right)=-b+1$, $f(\pi)=-a+b-c$,  $f\left(\dfrac{3\pi}{2}\right)=-b-1$.\\
		Suy ra $f(0)+f\left(\dfrac{\pi}{2}\right)+f(\pi)+f\left(\dfrac{3\pi}{2}\right)=0$ nên suy  ra nếu không có giá trị nào trong bốn giá trị bằng 0 thì ít nhất có một giá trị âm và dương.\\
		Từ đó suy ra  điều phải chứng minh.
	} 
\end{ex}

% nguồn: Ok_TuanLuan_Toan11_Lan2
\begin{ex}%[1D3V3-5]%[ma:650456f6]
	Chứng minh rằng phương trình $\sqrt{x^3+6x+1}-2=0$ có nghiệm dương. 
	 \loigiai{
		Ta có $\sqrt{x^3+6x+1}-2=0\Leftrightarrow \sqrt{x^3+6x+1}=2\Leftrightarrow x^3+6x+1=4\Leftrightarrow x^3+6x-3=0$.\\ 
		Xét hàm số $f(x)=x^3+6x-3$ là hàm số liên tục trên $\mathbb{R}$ nên nó liên tục trên $[0; 1]$.\\ 
		Ta có  $f(0)=-3$ và $f(1)=4$ nên $f(0) \cdot f(1) < 0$.\\
		Khi đó  phương trình $f(x)=0$ có ít nhất một nghiệm trong khoảng $(0; 1)$.\\
		Vậy phương trình $\sqrt{x^3+6x+1}-2=0$  có ít nhất một nghiệm dương.
	 }
\end{ex}

% nguồn: Ok_TuanLuan_Toan11_Lan2
\begin{ex}%[1D3V3-5]%[ma:7dfdadff]
	Chứng minh rằng phương trình $x^5-3x-7=0$ có ít nhất một  nghiệm.
	\loigiai{ 
		Xét hàm số $f(x)=x^5-3x-7$ là hàm đa thức liên tục trên $\mathbb{R}$ nên nó liên tục trên $[0;  2]$.\\
		Ta có $f(0)=-7, f(2)=19$ nên $f(0) \cdot f(2) < 0$. Khi đó phương trình $f(x)=0$ có ít nhất một nghiệm trong khoảng $(0;  2)$.
	} 
\end{ex}

% nguồn: Ok_TuanLuan_Toan11_Lan2
\begin{ex}%[1D3V3-5]%[ma:5318a5e8]
	Chứng minh rằng phương trình $a\cos 2x+b\sin x+\cos x=0$ luôn có nghiệm với mọi tham  số $a, b$.
	\loigiai{ 
		Đặt $f(x)=a\cos 2x+b\sin x+\cos x$ có tập xác định là $\mathscr{D}=\mathbb{R} \Rightarrow f(x)$  liên tục trên $\mathbb{R}$.\\
		$f(0)=a+1; f(\pi)=a-1; f\left(\dfrac{\pi}{2}\right)=-a+b;  f\left(\dfrac{3\pi}{2}\right)=-a-b$.\\
		Vì $f(0)+f(\pi)+f\left(\dfrac{\pi}{2}\right)+f\left(\dfrac{3\pi}{2}\right)=0$ nên trong bốn số phải có hai số mà tích của  chúng bé hơn hoặc bằng không.\\
		Vậy phương trình  $f(x)=0$ luôn có nghiệm với mọi tham số $a$, $b$.
	 }
\end{ex}

% nguồn: Ok_TuanLuan_Toan11_Lan2
\begin{ex}%[1D3V3-5]%[ma:efb8893c]
	Chứng minh rằng phương trình $6x^3+3x^2-31x+10=0$ có đúng  ba nghiệm phân biệt.
	 \loigiai{
		Đặt  $f(x)=6x^3+3x^2-31x+10$.\\
		 Tập xác định: $\mathscr{D}=\mathbb{R}$.\\
		$\Rightarrow f(x)$ liên tục trên $\mathbb{R} \Rightarrow f(x)$ liên tục  trên $[-3; 2]$.\\
		$\heva{&f(-3)=-32\\ &f(0)=10} \Rightarrow f(-3) \cdot  f(0) < 0\Rightarrow f(x)=0$ có nghiệm thuộc $(-3; 0)$.\\
		$\heva{&f(0)=10\\ &f(1)=-12} \Rightarrow f(0) \cdot f(1) <  0\Rightarrow f(x)=0$ có nghiệm thuộc $(1; 2)$.\\
		$\heva{&f(1)=-12\\ &f(2)=8}  \Rightarrow f(1) \cdot f(2) < 0\Rightarrow f(x)=0$ có nghiệm thuộc $(1; 2)$.\\
		Mặt khác vì $f(x)$ là một đa thức bậc ba nên $f(x)=0$ chỉ có tối đa ba nghiệm.\\ 
		Vậy phương trình $f(x)=0$  có đúng $3$ nghiệm phân biệt.
	 }
\end{ex}

% nguồn: Ok_TuanLuan_Toan11_Lan2
\begin{ex}%[1D3V3-5]%[ma:92f796c9]
	Chứng minh rằng phương  trình $x^3-2x^2+x-1-\dfrac{2}{x}+\dfrac{m}{x+3}=0$ luôn có nghiệm $\forall m$.
	\loigiai{ 
		Ta  có
		\[x^3-2x^2+x-1-\dfrac{2}{x}+\dfrac{m}{x+3}=0\Leftrightarrow  mx+(x-2)(x+3)(x^3+x+1)=0\ (*).\]
		Đặt  $f(x)=mx+(x-2)(x+3)(x^3+x+1)$.\\
		Nhận xét $f(x)$ liên tục trên  $\mathbb{R}$ và $\heva{&f(-3) \cdot f(0)=18m \\ &f(0) \cdot f(2)=-12m.}$\\
		\begin{itemize}
			\item \textbf{Trường hợp 1:} Nếu $m=0$, phương trình luôn có  1 nghiệm $x=2$.
			\item \textbf{Trường hợp  2:} Nếu $m \neq 0\Rightarrow (*)$ có ít nhất một nghiệm thuộc $(-3; 0)$ hoặc $(0; 2)$.
		\end{itemize}
		Suy ra $\forall  m \in \mathbb{R}$, luôn có ít nhất một nghiệm thuộc $\mathscr{D}$.\\
		Vậy $\forall m \in \mathbb{R}$, phương trình đã cho  có ít nhất một nghiệm.
	 }
\end{ex}
